\documentclass[10pt,a4paper]{amsart}
\usepackage{amsmath,amssymb,amsthm,graphicx}
\usepackage{verbatim}
\newenvironment{reference}{\comment}{\endcomment}
%\newenvironment{reference}{}{}
\newenvironment{slogan}{\comment}{\endcomment}
\newenvironment{history}{\comment}{\endcomment}
% For commutative diagrams we use Xy-pic
\usepackage[all]{xy}

% We use 2cell for 2-commutative diagrams.
\xyoption{2cell}
% Theorem environments.
%
\theoremstyle{plain}
\newtheorem{theorem}[subsection]{Theorem}
\newtheorem{proposition}[subsection]{Proposition}
\newtheorem{lemma}[subsection]{Lemma}

\theoremstyle{definition}
\newtheorem{definition}[subsection]{Definition}
\newtheorem{example}[subsection]{Example}
\newtheorem{exercise}[subsection]{Exercise}
\newtheorem{situation}[subsection]{Situation}

\theoremstyle{remark}
\newtheorem{remark}[subsection]{Remark}
\newtheorem{remarks}[subsection]{Remarks}

\numberwithin{equation}{subsection}

% Macros
%
\def\lim{\mathop{\mathrm{lim}}\nolimits}
\def\colim{\mathop{\mathrm{colim}}\nolimits}
\def\Spec{\mathop{\mathrm{Spec}}}
\def\Hom{\mathop{\mathrm{Hom}}\nolimits}
\def\Ext{\mathop{\mathrm{Ext}}\nolimits}
\def\SheafHom{\mathop{\mathcal{H}\!\mathit{om}}\nolimits}
\def\SheafExt{\mathop{\mathcal{E}\!\mathit{xt}}\nolimits}
\def\Sch{\mathit{Sch}}
\def\Mor{\mathop{\mathrm{Mor}}\nolimits}
\def\Ob{\mathop{\mathrm{Ob}}\nolimits}
\def\Sh{\mathop{\mathit{Sh}}\nolimits}
\def\NL{\mathop{N\!L}\nolimits}
\def\CH{\mathop{\mathrm{CH}}\nolimits}
\def\proetale{{pro\text{-}\acute{e}tale}}
\def\etale{{\acute{e}tale}}
\def\QCoh{\mathit{QCoh}}
\def\Ker{\mathop{\mathrm{Ker}}}
\def\Im{\mathop{\mathrm{Im}}}
\def\Coker{\mathop{\mathrm{Coker}}}
\def\Coim{\mathop{\mathrm{Coim}}}

% Boxtimes
%
\DeclareMathSymbol{\boxtimes}{\mathbin}{AMSa}{"02}

%
% Macros for moduli stacks/spaces
%
\def\QCohstack{\mathcal{QC}\!\mathit{oh}}
\def\Cohstack{\mathcal{C}\!\mathit{oh}}
\def\Spacesstack{\mathcal{S}\!\mathit{paces}}
\def\Quotfunctor{\mathrm{Quot}}
\def\Hilbfunctor{\mathrm{Hilb}}
\def\Curvesstack{\mathcal{C}\!\mathit{urves}}
\def\Polarizedstack{\mathcal{P}\!\mathit{olarized}}
\def\Complexesstack{\mathcal{C}\!\mathit{omplexes}}
% \Pic is the operator that assigns to X its picard group, usage \Pic(X)
% \Picardstack_{X/B} denotes the Picard stack of X over B
% \Picardfunctor_{X/B} denotes the Picard functor of X over B
\def\Pic{\mathop{\mathrm{Pic}}\nolimits}
\def\Picardstack{\mathcal{P}\!\mathit{ic}}
\def\Picardfunctor{\mathrm{Pic}}
\def\Deformationcategory{\mathcal{D}\!\mathit{ef}}
\usepackage{cleveref}
\crefname{originaldefinition}{original definition}{original definitions}
\Crefname{originaldefinition}{Original definition}{Original definitions}
\makeatletter
\newlabel{definition-P}{{definition-P}{source}{}{sourcepointer-1}{}}
\newlabel{definition-P@cref}{{[originaldefinition][0][]definition-P}{[1][0][]source}{}{}{}}
\makeatother
\crefname{originallemma}{original lemma}{original lemmas}
\Crefname{originallemma}{Original lemma}{Original lemmas}
\makeatletter
\newlabel{lemma-composition-finite-presentation}{{lemma-composition-finite-presentation}{source}{}{sourcepointer-2}{}}
\newlabel{lemma-composition-finite-presentation@cref}{{[originallemma][0][]lemma-composition-finite-presentation}{[1][0][]source}{}{}{}}
\makeatother
\crefname{originallemma}{original lemma}{original lemmas}
\Crefname{originallemma}{Original lemma}{Original lemmas}
\makeatletter
\newlabel{lemma-composition-flat}{{lemma-composition-flat}{source}{}{sourcepointer-3}{}}
\newlabel{lemma-composition-flat@cref}{{[originallemma][0][]lemma-composition-flat}{[1][0][]source}{}{}{}}
\makeatother
\crefname{originallemma}{original lemma}{original lemmas}
\Crefname{originallemma}{Original lemma}{Original lemmas}
\makeatletter
\newlabel{lemma-local-source-target}{{lemma-local-source-target}{source}{}{sourcepointer-4}{}}
\newlabel{lemma-local-source-target@cref}{{[originallemma][0][]lemma-local-source-target}{[1][0][]source}{}{}{}}
\makeatother
\crefname{originalreference}{original reference}{original references}
\Crefname{originalreference}{Original reference}{Original references}
\makeatletter
\newlabel{spaces-morphisms-lemma-smooth-flat}{{spaces-morphisms-lemma-smooth-flat}{source}{}{sourcepointer-5}{}}
\newlabel{spaces-morphisms-lemma-smooth-flat@cref}{{[originalreference][0][]spaces-morphisms-lemma-smooth-flat}{[1][0][]source}{}{}{}}
\makeatother
\crefname{originalreference}{original reference}{original references}
\Crefname{originalreference}{Original reference}{Original references}
\makeatletter
\newlabel{spaces-morphisms-lemma-smooth-locally-finite-presentation}{{spaces-morphisms-lemma-smooth-locally-finite-presentation}{source}{}{sourcepointer-6}{}}
\newlabel{spaces-morphisms-lemma-smooth-locally-finite-presentation@cref}{{[originalreference][0][]spaces-morphisms-lemma-smooth-locally-finite-presentation}{[1][0][]source}{}{}{}}
\makeatother
\crefname{originalreference}{original reference}{original references}
\Crefname{originalreference}{Original reference}{Original references}
\makeatletter
\newlabel{stacks-properties-lemma-composition-surjective}{{stacks-properties-lemma-composition-surjective}{source}{}{sourcepointer-7}{}}
\newlabel{stacks-properties-lemma-composition-surjective@cref}{{[originalreference][0][]stacks-properties-lemma-composition-surjective}{[1][0][]source}{}{}{}}
\makeatother
\title{A property of morphisms that is smooth local on source and target and fppf local on the target descends along a flat finitely presented surjection}
\author{The Stacks Project Authors}
\date{stacks-project:stacks-morphisms.tex:local-snapshot-20261012; GNU FDL 1.2 or later}
\begin{document}
\maketitle
\noindent\textit{Editorial note.} All delivered mathematical text is the author's own, copied from the pinned Stacks Project chapter source sources/stacks-algebraic-stacks-morphisms/pinned-original-source.tex; no statement or proof step was rewritten, reordered or removed.\par\medskip
\noindent\textit{Editorial note.} Preamble provenance: the Stacks Project keeps one single shared preamble for every chapter; this card takes it directly from the staged sources/stacks-algebraic-stacks-morphisms/preamble.tex as three verbatim spans (lines 8-12, 14-18 and 154-217), so no derived file was written and nothing was added or redefined. The wrapper class and the base packages are supplied by the assembler.\par\medskip
\noindent\textit{Editorial note.} Bibliography provenance: neither the counted statement nor its proof contains a citation (this lane rescans the whole span before writing the card), so the card carries no bibliography and no entry was reconstructed.\par\medskip
\noindent\textit{Editorial note.} Selection provenance: this card was opened under the file-based control contract (CONTROL.md), with the mechanical pre-card gate run on its statement and proof interval before the card was written.\par\medskip
\noindent\textit{Allowed standard prerequisites (editorial).}\begin{itemize}
\item The definition of a property of morphisms of algebraic stacks through presentations and its smooth locality on source and target. Invoked at the start of the counted proof (source lines 5791-5795). Turns the hypothesis on the base-changed morphism into a statement about presenting algebraic spaces.
\item The statements that a smooth morphism of algebraic spaces is flat and locally of finite presentation and that these properties are preserved by composition. Invoked inside the counted proof (source lines 5796-5805) as cross-chapter pointers. Makes the composite presentation morphism surjective flat and of finite presentation.
\item The smooth locality statement for the property on the source and target. Invoked inside the counted proof (source lines 5806-5816). Allows the test objects to be replaced by presentations.
\item The assumption that the property is fppf local on the target. Used at the end of the counted proof (source lines 5825-5829). Descends the property from the base-changed presenting morphism to the presenting morphism.
\end{itemize}\medskip
\ifdefined\hypertarget\hypertarget{sourcepointer-1}{}\fi
\noindent\textit{Source reference.} Original definition definition-P: sources/stacks-algebraic-stacks-morphisms/pinned-original-source.tex, source lines 3372--3372. Navigation pointer only: the counted text refers to the Stacks item definition-P, which is an existing labeled result of this same chapter and is not part of this card. The pointer carries the target's own label as its number and its definition line as its span; its content is not reproduced here.\par\medskip
\ifdefined\hypertarget\hypertarget{sourcepointer-2}{}\fi
\noindent\textit{Source reference.} Original lemma lemma-composition-finite-presentation: sources/stacks-algebraic-stacks-morphisms/pinned-original-source.tex, source lines 5614--5614. Navigation pointer only: the counted text refers to the Stacks item lemma-composition-finite-presentation, which is an existing labeled result of this same chapter and is not part of this card. The pointer carries the target's own label as its number and its definition line as its span; its content is not reproduced here.\par\medskip
\ifdefined\hypertarget\hypertarget{sourcepointer-3}{}\fi
\noindent\textit{Source reference.} Original lemma lemma-composition-flat: sources/stacks-algebraic-stacks-morphisms/pinned-original-source.tex, source lines 5363--5363. Navigation pointer only: the counted text refers to the Stacks item lemma-composition-flat, which is an existing labeled result of this same chapter and is not part of this card. The pointer carries the target's own label as its number and its definition line as its span; its content is not reproduced here.\par\medskip
\ifdefined\hypertarget\hypertarget{sourcepointer-4}{}\fi
\noindent\textit{Source reference.} Original lemma lemma-local-source-target: sources/stacks-algebraic-stacks-morphisms/pinned-original-source.tex, source lines 3237--3237. Navigation pointer only: the counted text refers to the Stacks item lemma-local-source-target, which is an existing labeled result of this same chapter and is not part of this card. The pointer carries the target's own label as its number and its definition line as its span; its content is not reproduced here.\par\medskip
\ifdefined\hypertarget\hypertarget{sourcepointer-5}{}\fi
\noindent\textit{Source reference.} Original reference spaces-morphisms-lemma-smooth-flat: sources/stacks-algebraic-stacks-morphisms/pinned-original-source.tex, source lines 5400--5400. Navigation pointer only: the counted text refers to the Stacks item spaces-morphisms-lemma-smooth-flat, which belongs to another chapter of the project and is not part of this card. The pointer carries the target's own label as its number because this chapter's staged snapshot does not contain the other chapter's tag table; its content is not reproduced here.\par\medskip
\ifdefined\hypertarget\hypertarget{sourcepointer-6}{}\fi
\noindent\textit{Source reference.} Original reference spaces-morphisms-lemma-smooth-locally-finite-presentation: sources/stacks-algebraic-stacks-morphisms/pinned-original-source.tex, source lines 304--304. Navigation pointer only: the counted text refers to the Stacks item spaces-morphisms-lemma-smooth-locally-finite-presentation, which belongs to another chapter of the project and is not part of this card. The pointer carries the target's own label as its number because this chapter's staged snapshot does not contain the other chapter's tag table; its content is not reproduced here.\par\medskip
\ifdefined\hypertarget\hypertarget{sourcepointer-7}{}\fi
\noindent\textit{Source reference.} Original reference stacks-properties-lemma-composition-surjective: sources/stacks-algebraic-stacks-morphisms/pinned-original-source.tex, source lines 5403--5403. Navigation pointer only: the counted text refers to the Stacks item stacks-properties-lemma-composition-surjective, which belongs to another chapter of the project and is not part of this card. The pointer carries the target's own label as its number because this chapter's staged snapshot does not contain the other chapter's tag table; its content is not reproduced here.\par\medskip
\begin{lemma}
\label{lemma-descent-property}
Let $\mathcal{P}$ be a property of morphisms of algebraic spaces
which is smooth local on the source-and-target and fppf local
on the target.
Let $f : \mathcal{X} \to \mathcal{Y}$ be a morphism of algebraic stacks.
Let $\mathcal{Z} \to \mathcal{Y}$ be a surjective, flat, locally finitely
presented morphism of algebraic stacks. If the base change
$\mathcal{Z} \times_\mathcal{Y} \mathcal{X} \to \mathcal{Z}$
has $\mathcal{P}$, then $f$ has $\mathcal{P}$.
\end{lemma}

\begin{proof}
Assume $\mathcal{Z} \times_\mathcal{Y} \mathcal{X} \to \mathcal{Z}$
has $\mathcal{P}$. Choose an algebraic space $W$ and a surjective
smooth morphism $W \to \mathcal{Z}$. Observe that
$W \times_\mathcal{Z} \mathcal{Z} \times_\mathcal{Y} \mathcal{X} =
W \times_\mathcal{Y} \mathcal{X}$. Thus by the very definition of
what it means for $\mathcal{Z} \times_\mathcal{Y} \mathcal{X} \to \mathcal{Z}$
to have $\mathcal{P}$ (see Definition \ref{definition-P}
and Lemma \ref{lemma-local-source-target})
we see that $W \times_\mathcal{Y} \mathcal{X} \to W$
has $\mathcal{P}$. On the other hand, $W \to \mathcal{Z}$
is surjective, flat, and locally of finite presentation
(Morphisms of Spaces, Lemmas
\ref{spaces-morphisms-lemma-smooth-flat} and
\ref{spaces-morphisms-lemma-smooth-locally-finite-presentation})
hence $W \to \mathcal{Y}$ is surjective, flat, and locally of finite
presentation (by
Properties of Stacks, Lemma
\ref{stacks-properties-lemma-composition-surjective}
and
Lemmas \ref{lemma-composition-flat} and
\ref{lemma-composition-finite-presentation}).
Thus we may replace $\mathcal{Z}$ by $W$.

\medskip\noindent
Choose an algebraic space $V$ and a surjective smooth morphism
$V \to \mathcal{Y}$. Choose an algebraic space $U$ and a surjective
smooth morphism $U \to V \times_\mathcal{Y} \mathcal{X}$.
We have to show that $U \to V$ has $\mathcal{P}$.
Now we base change everything by $W \to \mathcal{Y}$: Set
$U' = W \times_\mathcal{Y} U$,
$V' = W \times_\mathcal{Y} V$,
$\mathcal{X}' = W \times_\mathcal{Y} \mathcal{X}$,
and $\mathcal{Y}' = W \times_\mathcal{Y} \mathcal{Y} = W$.
Then it is still true that $U' \to V' \times_{\mathcal{Y}'} \mathcal{X}'$
is smooth by base change. Hence by Lemma \ref{lemma-local-source-target}
used in the definition of $\mathcal{X}' \to \mathcal{Y}' = W$
having $\mathcal{P}$ we see that $U' \to V'$ has $\mathcal{P}$.
Then, since $V' \to V$ is surjective, flat, and locally of finite presentation
as a base change of $W \to \mathcal{Y}$ we see that $U \to V$
has $\mathcal{P}$ as $\mathcal{P}$ is local in the fppf topology
on the target.
\end{proof}

\end{document}
